\begin{frame}[fragile]{child\_process JS Library}

\begin{itemize}
	\item We can use Jest to run all tests in test directories. 
	\item But what if we need to use other tests than Jest?
	\item In this case, we should execute the tests one by one in an application. 
	\item JavaScript child\_process library is the tool for this. 
\end{itemize}

\end{frame}

\begin{frame}[fragile]{}

\begin{itemize}
	\item This is an example that uses the exec() method in the child\_process library. 
	\item We can give any command (line 3), and check the results or error conditions (lines 4 -- 9). 
\end{itemize}

\begin{Code}{}%
\begin{lstlisting}[firstnumber=1, label=glabels, xleftmargin=0pt, basicstyle=\scriptsize]% 

var exec = require('child_process').exec;

exec('cat *.js | wc -l',
    function (error, stdout, stderr) {
        console.log('stdout: ' + stdout);
        console.log('stderr: ' + stderr);
        if (error !== null) {
             console.log('exec error: ' + error);
        }
    });
\end{lstlisting}%   
\end{Code}%
\end{frame}

\begin{frame}[fragile]{Regression Test with Exec()}

\begin{itemize}
	\item Using the exec() method, we can run all the tests to check if any change breaks the existing test. 
\end{itemize}

\begin{Code}{}%
\begin{lstlisting}[firstnumber=1, label=glabels, xleftmargin=0pt, basicstyle=\scriptsize]% 

...
exec('jest',
    function (error, stdout, stderr) { ... });
exec('gui_test1',
    function (error, stdout, stderr) { ... });    
exec('gui_test2',
    function (error, stdout, stderr) { ... });    
\end{lstlisting}%   
\end{Code}%

\end{frame}